% 第四章 基于RDHO的任务卸载策略 - 公式汇总
% LaTeX版本

\documentclass[12pt]{article}
\usepackage{amsmath}
\usepackage{amssymb}
\usepackage{ctex}

\begin{document}

\title{第四章公式汇总 - LaTeX版本}
\maketitle

%% ==================== 4.1 系统模型与问题描述 ====================

\section{4.1 系统模型与问题描述}

\subsection{4.1.1 用户体验质量（QoE）模型}

% 公式4-1：时延满足度
\textbf{公式4-1：时延满足度}
\begin{equation}
S_i^{delay} = \frac{1}{1 + e^{\gamma_i \cdot (T_i - T_i^{max})}}
\end{equation}

% 公式4-2：能耗接受度
\textbf{公式4-2：能耗接受度}
\begin{equation}
S_i^{energy} = \begin{cases}
1 - \dfrac{E_i}{E_i^{budget}}, & \text{if } B_i \geq 60\% \\[10pt]
1 - \dfrac{E_i}{E_i^{budget}} \cdot (1 + \eta \cdot (0.6 - B_i)), & \text{if } 30\% \leq B_i < 60\% \\[10pt]
1 - \dfrac{E_i}{E_i^{budget}} \cdot (1 + 2\eta \cdot (0.3 - B_i)), & \text{if } B_i < 30\%
\end{cases}
\end{equation}

% 公式4-3：任务完成率
\textbf{公式4-3：任务完成率}
\begin{equation}
S_i^{completion} = \begin{cases}
1, & \text{if } T_i \leq T_i^{max} \text{ and } AoI_i \leq AoI_i^{max} \\
0, & \text{otherwise}
\end{cases}
\end{equation}

% 公式4-4：综合QoE评分
\textbf{公式4-4：综合QoE评分}
\begin{equation}
QoE_i = w_1^{qoe} \cdot S_i^{delay} + w_2^{qoe} \cdot S_i^{energy} + w_3^{qoe} \cdot S_i^{completion}
\end{equation}

% 公式4-5：系统平均QoE
\textbf{公式4-5：系统平均QoE}
\begin{equation}
\overline{QoE} = \frac{1}{K} \sum_{i=1}^{K} QoE_i
\end{equation}

\subsection{4.1.2 公平性度量与建模}

% 公式4-6：Jain公平性指数
\textbf{公式4-6：Jain公平性指数}
\begin{equation}
Fairness = \frac{\left(\sum_{i=1}^{K} QoE_i\right)^2}{K \cdot \sum_{i=1}^{K} QoE_i^2}
\end{equation}

\subsection{4.1.3 优化问题建模}

% 公式4-7：目标函数
\textbf{公式4-7：目标函数}
\begin{equation}
\min_{X} F(X) = \sum_{i \in \mathcal{T}} \left( w_E \cdot \frac{E_i}{E^{norm}} + w_T \cdot \frac{T_i}{T^{norm}} + w_{AoI} \cdot \frac{\overline{AoI}}{AoI^{norm}} - w_{QoE} \cdot QoE_i - w_F \cdot Fairness \right)
\end{equation}

% 公式4-8：任务分配唯一性约束
\textbf{公式4-8：任务分配唯一性约束}
\begin{equation}
x_i^{local} + x_i^{edge} + x_i^{cloud} = 1, \quad \forall i \in \mathcal{T}
\end{equation}

% 公式4-9：资源容量约束
\textbf{公式4-9：资源容量约束}
\begin{equation}
\sum_{i \in \mathcal{T}} x_i \cdot f_{i,j} \leq F_j^{max}, \quad \forall j \in \mathcal{E} \cup \mathcal{C}
\end{equation}

% 公式4-10：时延约束
\textbf{公式4-10：时延约束}
\begin{equation}
T_i \leq T_i^{max}, \quad \forall i \in \mathcal{T}
\end{equation}

% 公式4-11：AoI约束
\textbf{公式4-11：AoI约束}
\begin{equation}
AoI_i \leq AoI_i^{max}, \quad \forall i \in \mathcal{T}
\end{equation}

% 公式4-12：服务质量约束
\textbf{公式4-12：服务质量约束}
\begin{equation}
\frac{1}{K} \sum_{i=1}^{K} \mathbb{I}(T_i \leq T_i^{max}) \geq 0.95
\end{equation}

% 公式4-13：频率范围约束
\textbf{公式4-13：频率范围约束}
\begin{equation}
f_j^{min} \leq f_{i,j} \leq f_j^{max}, \quad \forall i \in \mathcal{T}, j \in \mathcal{D} \cup \mathcal{E} \cup \mathcal{C}
\end{equation}

% 公式4-14：QoE约束
\textbf{公式4-14：QoE约束}
\begin{equation}
QoE_i \geq QoE_{min}, \quad \forall i \in \mathcal{T}
\end{equation}

% 公式4-15：公平性约束
\textbf{公式4-15：公平性约束}
\begin{equation}
Fairness \geq 0.8
\end{equation}

%% ==================== 4.2 RDHO混合元启发式算法设计 ====================

\section{4.2 RDHO混合元启发式算法设计}

\subsection{4.2.2 霜冰优化算法（RIME）简介}

% 公式4-16：软凝华搜索
\textbf{公式4-16：软凝华搜索}
\begin{equation}
X_{new} = X_{best} + \beta \cdot \cos(\theta) \cdot h \cdot (ub - lb)
\end{equation}

% 公式4-17：自适应搜索半径
\textbf{公式4-17：自适应搜索半径}
\begin{equation}
h = 2 \left(1 - \frac{t}{T_{max}}\right)
\end{equation}

% 公式4-18：硬凝华穿刺
\textbf{公式4-18：硬凝华穿刺}
\begin{equation}
X_{new}[j_{rand}] = X_{best}[j_{rand}]
\end{equation}

% 公式4-19：穿刺概率
\textbf{公式4-19：穿刺概率}
\begin{equation}
E = 2 \cdot \exp\left(-\left(\frac{4t}{T_{max}}\right)^2\right)
\end{equation}

\subsection{4.2.3 蜣螂优化算法（DBO）简介}

% 公式4-20：滚球行为
\textbf{公式4-20：滚球行为}
\begin{equation}
X_{new} = X + \alpha \cdot k \cdot (X - lb) + b \cdot |X - X_{worst}|
\end{equation}

% 公式4-21：滚球递减因子
\textbf{公式4-21：滚球递减因子}
\begin{equation}
\alpha = 1 - \frac{t}{T_{max}}
\end{equation}

% 公式4-22：繁殖行为
\textbf{公式4-22：繁殖行为}
\begin{equation}
X_{new} = X_{best} + \beta_1 \cdot |X - X_{local\_best}|
\end{equation}

% 公式4-23：觅食行为
\textbf{公式4-23：觅食行为}
\begin{equation}
X_{new} = X_{best} + C_1 \cdot (X - lb) + C_2 \cdot (X - ub)
\end{equation}

% 公式4-24：偷窃行为
\textbf{公式4-24：偷窃行为}
\begin{equation}
X_{new} = X_{local\_best} + \tan(\theta) \cdot |X - X_{local\_best}|
\end{equation}

\subsection{4.2.4 RDHO融合机制}

% 公式4-25：双源混合初始化
\textbf{公式4-25：双源混合初始化}
\begin{equation}
X_i = \begin{cases}
\mathcal{N}(\mu, \sigma^2), & \text{前50\%个体（高斯分布）} \\
\mathcal{U}(lb, ub), & \text{后50\%个体（均匀分布）}
\end{cases}
\end{equation}

% 公式4-26：生产者融合更新
\textbf{公式4-26：生产者融合更新}
\begin{equation}
X_{new} = w \cdot X_{RIME} + (1 - w) \cdot X_{DBO}
\end{equation}

% 公式4-27：自适应融合权重
\textbf{公式4-27：自适应融合权重}
\begin{equation}
w = 0.5 + 0.3 \cdot \cos\left(\frac{\pi \cdot t}{T_{max}}\right)
\end{equation}

% 公式4-28：跟随者条件切换
\textbf{公式4-28：跟随者条件切换}
\begin{equation}
X_{new} = \begin{cases}
\text{RIME硬凝华穿刺}, & \text{if } rand < E \\
\text{DBO觅食}, & \text{otherwise}
\end{cases}
\end{equation}

% 公式4-29：侦察者更新
\textbf{公式4-29：侦察者更新}
\begin{equation}
X_{new} = \begin{cases}
\text{DBO偷窃}, & \text{if } f(X) > f(X_{best}) \\
X_{best} + scale \cdot X_{best} \cdot Cauchy(), & \text{otherwise}
\end{cases}
\end{equation}

% 公式4-30：Cauchy变异尺度
\textbf{公式4-30：Cauchy变异尺度}
\begin{equation}
scale = 0.1 \cdot \left(1 - \frac{t}{T_{max}}\right)
\end{equation}

% 公式4-31：动态分层惩罚
\textbf{公式4-31：动态分层惩罚}
\begin{equation}
penalty = base\_penalty \cdot (1 + 2 \cdot progress)^{\alpha}
\end{equation}

% 公式4-32：搜索进度
\textbf{公式4-32：搜索进度}
\begin{equation}
progress = \frac{t}{T_{max}}
\end{equation}

\subsection{4.2.6 算法时间复杂度分析}

% 公式4-33：时间复杂度
\textbf{公式4-33：时间复杂度}
\begin{equation}
O(T_{max} \cdot N \cdot (K + \log N))
\end{equation}

\textbf{demo}
\begin{equation}
E = 2 \cdot \exp\left(-\left(\frac{4t}{T_{max}}\right)^2\right)
\end{equation}

\end{document}
